% Ce chapitre repose sur :
%   DOC: s30
%   OBS: REL-DEM.P05
%   HYP: (aucun)
%   CHI: dms-affichee, dms-publie, exigences-calculables, exigences-champs-manquants, exigences-classement-mixte, exigences-ecart-maximal
%   CHI: exigences-non-calculables, exigences-rubriques, exigences-sans-objet, facteur-annualisation, irot-affiche, irot-publie
%   CHI: rotation-affichee, rotation-affichee-avant, tdb-filtrabilite-non, tdb-filtrabilite-oui, tdb-filtrabilite-sous-reserve, tdb-graphiques
%   CHI: tdb-indicateurs, tdb-indicateurs-methode, tdb-indicateurs-pilotage, tdb-page-plus-chargee, tdb-pages, tdb-pages-methode
%   CHI: tdb-pages-pilotage, tdb-recalcules-depuis-faits, tolerance-indicateurs-sejour, tom-affiche, tom-publie, trot-publie
%   SER: exigences-restitution, tableau-de-bord-par-page
%   SECTIONS: 3

\chapter{Le tableau de bord}\label{chap:tableau-de-bord}

Le tableau de bord donne au responsable du service d'accueil et d'admission les grandeurs de son
activité, calculées chaque jour depuis les données de la chaîne. Chaque indicateur y affiche sa
définition et la décision qu'il sert. La question traitée est la manière dont le responsable s'en
sert, et ce que le tableau de bord couvre des exigences de la statistique nationale.

\section{Conception}

Le tableau de bord s'adresse au responsable du service, et non à un analyste. Chaque indicateur
affiche sa définition en une phrase, la décision qu'elle sert et la source de sa valeur.
Sur les \chiffre{tdb-indicateurs} indicateurs, \chiffre{tdb-recalcules-depuis-faits} sont
recalculés depuis les faits, comme le cadrage le demande. Les autres sont des écarts assumés, que l'écran
nomme.

Les pages lisent un schéma dédié de tables figées, daté, que la chaîne reconstruit puis remplace
d'un bloc en une transaction. Ce choix évite au lecteur une erreur de lecture pendant la
reconstruction, et l'annexe~\ref{annexe:complements} en donne la mesure
(section~\ref{sec:complements-instantane}).

Les pages se répartissent en deux sections nommées par leur public. Un indicateur reste en
« Pilotage du service » si sa valeur peut changer une décision du service. Il passe en « Évaluation
de la chaîne » s'il décrit la chaîne, sa performance, sa provenance ou ses paramètres. La règle
donne \chiffre{tdb-indicateurs-pilotage} indicateurs en pilotage et
\chiffre{tdb-indicateurs-methode} en évaluation. Les indicateurs que le registre des relations
injectées déclare circulaires sont tous en évaluation, et un contrôle le vérifie.

Le tableau de bord compte \chiffre{tdb-pages} pages, dont \chiffre{tdb-pages-pilotage} en pilotage
et \chiffre{tdb-pages-methode} en évaluation. Elles portent \chiffre{tdb-indicateurs} indicateurs et
\chiffre{tdb-graphiques} tracés (cf. tableau~\ref{tab:pages-tdb}).

\begin{table}[htbp]
\centering
\caption[Les pages du tableau de bord]{Les pages du tableau de bord, dans l'ordre du registre des indicateurs}\label{tab:pages-tdb}
\footnotesize
\pgfplotstabletypeset[
  col sep=comma,
  % LE NOM DE PAGE SE COMPOSE EN CHASSE FIXE, comme le nom de table au chapitre précédent et comme
  % tout identifiant du document. En romain, `\textunderscore` produit une barre basse démesurée :
  % `rendez_vous` s'y lisait `rendez__vous`.
  columns/page/.style={column name={Page}, string type,
    string replace*={_}{\_},
    postproc cell content/.append style={@cell content/.add={\texttt\bgroup}{\egroup}}},
  columns/section/.style={column name={Section}, string type},
  columns/indicateurs/.style={column name={Indicateurs}, int detect},
  columns/graphiques/.style={column name={Tracés}, int detect},
  every head row/.style={before row=\toprule, after row=\midrule},
  every last row/.style={after row=\bottomrule},
]{\serie{tableau-de-bord-par-page}}
\end{table}

Chaque page y porte sa section, le nombre d'indicateurs qu'elle affiche et le nombre de tracés que
son fichier appelle. En pilotage, chaque page sert une décision : les plages d'ouverture pour
\emph{Activité}, les créneaux pour \emph{Rendez-vous}, l'accueil et le tri pour \emph{Urgences}, les
lits pour \emph{Séjours}, les relances pour \emph{Facturation et recouvrement}, la plus chargée avec
\chiffre{tdb-page-plus-chargee} indicateurs, la saisie pour
\emph{Qualité des données}, et l'extraction de lignes pour \emph{Données}.

\section{Un scénario d'utilisation}

Le responsable du service se demande s'il faut ouvrir des créneaux de rendez-vous dans une activité.
Trois pages lui suffisent pour y répondre.

Il ouvre d'abord la page \emph{Activité}, qui s'affiche par défaut (voir
figure~\ref{fig:tdb-activite}). La barre latérale y montre les deux sections, et le haut de la page
l'état de l'instantané et le filtre de période. Le volume d'événements par jour, ventilé par
famille, lui dit si l'activité programmée monte ou baisse sur la période choisie.

\begin{figure}[htbp]
\centering
\captureTdb{page-activite}{0.92\linewidth}
\caption[La page Activité du tableau de bord]{La page \emph{Activité}, première de la section de pilotage : le volume d'événements par jour, ventilé par famille}\label{fig:tdb-activite}
\end{figure}

Il passe ensuite à la page \emph{Rendez-vous}. Il y lit le délai d'obtention en médiane et au
neuvième décile, le taux d'absence et le taux d'annulation, par activité. La page sépare
l'annulation, qui laisse réaffecter le créneau, de l'absence, qui le perd. Chaque bloc dit aussi si
l'indicateur suit le filtre de période. Au total, \chiffre{tdb-filtrabilite-oui} indicateurs le
suivent sans réserve, \chiffre{tdb-filtrabilite-sous-reserve} par une date autre que celle qu'un lecteur
attendrait, et \chiffre{tdb-filtrabilite-non} ne le suivent pas. L'écran nomme alors le motif.

Il en tire une première réponse. Un délai qui s'allonge avec peu d'absences appelle des créneaux en
plus. Des absences nombreuses appellent d'abord un rappel des patients, puisque chaque absence perd
un créneau.

Pour appuyer sa décision, il ouvre enfin la page \emph{Données} (voir figure~\ref{fig:tdb-donnees}).
Il y filtre les rendez-vous par service et par activité, puis télécharge la sélection en un fichier
tabulaire. Le fichier suit le format des livrables de la chaîne : séparateur virgule, encodage
universel avec marque d'ordre d'octets. Un tableur l'ouvre par double-clic sans abîmer les accents.

\begin{figure}[htbp]
\centering
\captureTdb{page-donnees}{0.92\linewidth}
\caption[La page Données du tableau de bord]{La page \emph{Données} : les trois filtres, le sélecteur de table, le décompte des lignes retenues et le bouton de téléchargement}\label{fig:tdb-donnees}
\end{figure}

\section{Les exigences de la statistique nationale}

Le tableau de bord couvre la plupart des rubriques que la statistique nationale attend, et nomme le
champ qui bloque les autres. Le logiciel observé porte un profil dédié aux rapports et aux
statistiques\releve{REL-DEM.P05}. Ce que ce profil doit produire se lit dans la nomenclature
nationale \cite{s30}, dont quatre tableaux seulement descendent à l'établissement nommé. Ces quatre
tableaux donnent la liste des rubriques attendues (cf. tableau~\ref{tab:exigences-restitution}).

\begin{table}[htbp]
\centering
\caption[Les rubriques attendues par la nomenclature nationale]{Les rubriques que la nomenclature nationale attend d'un établissement nommé, et l'état de chacune}\label{tab:exigences-restitution}
\footnotesize
\pgfplotstabletypeset[
  col sep=tab, string type,
  columns={indicateur,classement},
  columns/indicateur/.style={column name={Rubrique attendue}, column type={p{6.4cm}}},
  columns/classement/.style={column name={État}, column type={p{6.4cm}}},
  every head row/.style={before row=\toprule, after row=\midrule},
  every last row/.style={after row=\bottomrule},
]{\serie{exigences-restitution}}
\end{table}

Les rubriques y sont au nombre de \chiffre{exigences-rubriques}, dont
\chiffre{exigences-calculables} calculables, \chiffre{exigences-non-calculables} non calculables,
\chiffre{exigences-sans-objet} sans objet et \chiffre{exigences-classement-mixte} au classement
mixte, celle des laboratoires.
Une rubrique non calculable faute de champ désigne un manque du système d'information. Une rubrique
sans objet désigne une activité que l'établissement n'exerce pas, et rien n'y est à reprocher au
logiciel. Les tables mobilisées et le champ manquant de chaque rubrique figurent au registre des
indicateurs.

Les champs manquants sont au nombre de \chiffre{exigences-champs-manquants}, et chacun bloque une
seule rubrique. Le référentiel des médecins vient en tête : son absence produit un écart de
\chiffre{exigences-ecart-maximal}~\% sur un indicateur dont le numérateur se calcule. Les deux
autres champs rendent leur rubrique simplement absente.

La page \emph{Séjours} affiche les quatre indicateurs réglementaires annualisés, chacun avec la
source de sa valeur (voir figure~\ref{fig:tdb-sejours}).

\begin{figure}[htbp]
\centering
\captureTdb{page-sejours}{0.92\linewidth}
\caption[La page Séjours du tableau de bord]{La page \emph{Séjours} : les indicateurs réglementaires annualisés, chacun affichant la source de sa valeur}\label{fig:tdb-sejours}
\end{figure}

L'un d'eux a été longtemps faux à l'écran. Le taux de rotation valait
\chiffre{rotation-affichee-avant} là où la statistique nationale publie \chiffre{trot-publie}. L'écart
valait exactement \chiffre{facteur-annualisation}, le rapport de la durée de la période à l'année :
la page n'annualisait pas la grandeur, alors que le registre la déclarait annualisée. Les contrôles
de l'époque interrogeaient la couche analytique, et celui de la page recopiait sa formule pour se
comparer à elle. Il est resté vert pendant toute la durée du défaut.

Le contrôle qui l'a remplacé appelle la page et confronte ses sorties aux valeurs publiées et à une
implémentation écrite séparément dans un autre langage. Les quatre valeurs affichées valent
aujourd'hui \chiffre{tom-affiche}~\% d'occupation contre \chiffre{tom-publie}~\% publiés,
\chiffre{dms-affichee} jours de durée moyenne de séjour contre \chiffre{dms-publie},
\chiffre{rotation-affichee} de rotation contre \chiffre{trot-publie}, et \chiffre{irot-affiche} jours
d'intervalle contre \chiffre{irot-publie}. Les quatre restent sous la tolérance relative de
\chiffre{tolerance-indicateurs-sejour} que le contrôle porte.

\section*{Synthèse}

Le tableau de bord met sous les yeux du responsable les grandeurs de son service, avec leur
définition et la décision qu'elles servent. Un responsable passe de la page \emph{Activité} à la
page \emph{Rendez-vous}, puis aux \emph{Données}, pour décider d'ouvrir des créneaux. Le tableau de
bord calcule \chiffre{exigences-calculables} des \chiffre{exigences-rubriques} rubriques attendues
par la statistique nationale, et nomme le champ qui manque aux autres. Le
chapitre~\ref{chap:recommandations} en tire des recommandations pour le service.
